\documentclass[11pt]{article}
\usepackage{fontspec}
\setmainfont{Times New Roman}
\setsansfont{Arial}
\setmonofont{Consolas}
\usepackage{amsmath,amssymb,mathtools}
\usepackage{geometry}
\usepackage{microtype}
\geometry{a4paper,margin=2.28cm}
\setlength{\parindent}{1.15em}
\setlength{\parskip}{0pt}
\providecommand{\mZ}{\mathfrak{Z}}
\emergencystretch=120em
\allowdisplaybreaks
\sloppy
\hbadness=10000
\vbadness=10000
\hfuzz=2pt
\vfuzz=10pt
\raggedbottom
\begin{document}
% Унит: Папер34 Сецтион22 соурце-фиделиты цонтинуатион в001.
% Лангуаге лане: Интерславиц.
% Дирецт соурце: соурцес/папер34/соурце_фиделиты/Noether_Папер34_Сецтион22_ОРИГИНАЛ_СЦАН_ВИТНЕСС_в001.текс.
\section*{Emmy Noether: Хиперкомплексне величины и теорија представјень}
\emph{Папир 34, §22. Изворно-точны превод из сканованого немецкого сведка.}

\subsection*{§ 22. Применьенье к абеловым групам}

Нехај
\[
        G=G_1\times G_2\times\cdots\times G_r
\]
буде конечна Абелова група, разложена в директны продукт цикличных груп $G_i$ порјадков $h_i$. Јеј елементы сут
\[
        a=a_1^{\alpha_1}\cdots a_r^{\alpha_r}.
\]
Нехај $P$ буде тело, чија карактеристика не дели порјадок групы
\[
        h=h_1h_2\cdots h_r.
\]

Групово колцо $\mZ$ состоји из всех сум
\[
        \sum_{\alpha_1,\ldots,\alpha_r}
        A_{\alpha_1\ldots\alpha_r}a_1^{\alpha_1}\cdots a_r^{\alpha_r},
        \qquad A_{\alpha_1\ldots\alpha_r}\in P,
\]
и јест хомоморфны образ полиномиалного колца $P[x_1,\ldots,x_r]$ посредством $x_i\mapsto a_i$. Значи
\[
        \mZ\simeq P[x_1,\ldots,x_r]/\mathfrak m,
        \qquad
        \mathfrak m=(x_1^{h_1}-e,\ldots,x_r^{h_r}-e).
\]
В разширном телу $\Omega$ полиномы $x_i^{h_i}-e$ разпадајут се на саме различне факторы $x_i-\xi_i$; зато $\mathfrak m$ става продуктом (или пресеком) самых различних простых идеалов
\[
        (x_1-\xi_1^{(\alpha_1)},\ldots,x_r-\xi_r^{(\alpha_r)}),
\]
всакому из ньих одповедаје једна нульева точка
\[
        (\xi_1^{(\alpha_1)},\ldots,\xi_r^{(\alpha_r)}).
\]
Твореньу пресека одповедаје директно сумно разложенье (§ 4):
\[
        \mZ\Omega=\mZ_1+\cdots+\mZ_h,
\]
кде всакы раз $\mZ_i\sim\Omega$; зато достајемо представјенье
\[
        a_i\mapsto \xi_i^{(\alpha_i)},
        \qquad\hbox{или обче}\qquad
        a\mapsto \chi^{(\alpha)}(a).
\]
Та представјеньа сут карактере. Ньихово количство јест
\[
        h_1h_2\cdots h_r=h,
\]
како то требује обча теорија.

\end{document}

